\documentclass[11pt]{article}

\usepackage[a4paper,margin=27mm]{geometry}
\usepackage{amsmath,amssymb,amsthm,mathtools}
\usepackage[hidelinks]{hyperref}
\usepackage{microtype}
\setlength{\emergencystretch}{3em}

\newtheorem{theorem}{Theorem}[section]
\newtheorem{lemma}[theorem]{Lemma}
\newtheorem{proposition}[theorem]{Proposition}
\newtheorem{corollary}[theorem]{Corollary}
\theoremstyle{definition}
\newtheorem{definition}[theorem]{Definition}
\theoremstyle{remark}
\newtheorem{remark}[theorem]{Remark}

\newcommand{\Z}{\mathbb Z}
\newcommand{\Rel}{\mathcal R}
\newcommand{\Sym}{\operatorname{Sym}}
\newcommand{\ev}{\operatorname{ev}}
\newcommand{\src}{\operatorname{src}}
\newcommand{\im}{\operatorname{im}}

\title{Implementation-Readiness Audit of the\
Exceptional Terminal PBW Lifting Argument}
\author{}
\date{August 2026}

\begin{document}
\maketitle

\begin{abstract}
This note is an erratum and implementation audit of the proposed evaluated
cokernel--Stokes proof of the exceptional terminal lifting lemma.  A second
check against the exact Lean objects reveals a fatal omission: the external
boundary of the complete rewrite trace contains its initial root family.
Its primitive read is the governing PBW primitive, not a certified zero.
Once that term is restored, the proposed terminal-package vanishing is no
longer available.

The failure is not merely a matter of exposition.  The current Lean library
proves, by a short calculation from already verified aggregate identities,
that the canonical exceptional discrepancy evaluates to the negative of the
governing element.  Thus proving that discrepancy zero is equivalent to the
conclusion of Step~7; it cannot serve as an independent local lemma.  The
calculation was checked by compiling the precise Lean theorem stated in
Section~5 below.

The evaluated-range formulation and its lifting criterion remain correct,
but they do not prove range membership.  Consequently the earlier proof is
withdrawn and should not be implemented.  The end of the note gives the
correct root-inclusive target and a four-point implementation plan which
does not hide the missing assertion behind a quotient or a new terminal
package.
\end{abstract}

\section{Verdict}

The proposed proof is not implementation-ready and is not a valid proof of
the exceptional lifting lemma.  Its finite Stokes calculation used the
identity
\[
                         I-F=\partial_{\mathrm{led}}S
\]
for a complete labelled rewrite forest, but its claimed external-boundary
partition omitted the initial family \(I\).  The omission changes the
conclusion from ``the exceptional obstruction is zero'' to ``the
exceptional obstruction equals the remaining root obstruction.''  That root
obstruction is precisely the governing top-layer element which Step~7 must
annihilate.

There is therefore no safe short implementation of the earlier document.
In particular, introducing an evaluated cokernel, a free occurrence group,
or a new terminal-package collector would add infrastructure without
removing the mathematical obstruction.

\section{The lifting statement and the part which remains correct}

Let \(F=\bigoplus_{s=1}^{n+1}F_s\) be the truncated free metabelian Lie
ring, let \(\ev:F\to L\) be evaluation, and put
\(\Rel=\ker(\ev)\).  Set
\[
 A=\bigoplus_{s=1}^{n}F_s,
 \qquad D=\im(\Rel\to A),
 \qquad K=D\otimes A.
\]
The terminal Koszul differential and complete source primitive are additive
maps
\[
 \partial:K\longrightarrow\Sym^2A,
 \qquad
 \src:K\longrightarrow F.
\]
For the exceptional signed relation-left family write
\(B_{\mathcal E}\in\Sym^2A\) for its complete factor-two read and
\(P_{\mathcal E}\in F\) for its complete factor-one PBW read.  The desired
lemma asks for \(z_{\mathcal E}\in K\) and \(r_{\mathcal E}\in\Rel\) such
that
\begin{equation}
 \partial z_{\mathcal E}=B_{\mathcal E},
 \qquad
 \src z_{\mathcal E}=P_{\mathcal E}+r_{\mathcal E}.       \label{eq:target}
\end{equation}

The simultaneous range formulation is correct.  Define
\[
 \Lambda_{\mathrm{ev}}:K\longrightarrow\Sym^2A\oplus L,
 \qquad
 \Lambda_{\mathrm{ev}}(z)=(\partial z,\ev(\src z)).
\]

\begin{lemma}[simultaneous lifting criterion]\label{lem:range}
For \(b\in\Sym^2A\) and \(p\in F\), the pair
\((b,\ev p)\) belongs to \(\im\Lambda_{\mathrm{ev}}\) if and only if there
are \(z\in K\) and \(r\in\Rel\) such that
\[
                         \partial z=b,
 \qquad                  \src z=p+r.
\]
\end{lemma}

\begin{proof}
Membership in the range gives \(z\) with
\(\partial z=b\) and \(\ev(\src z-p)=0\).  Put
\(r=\src z-p\).  Then \(r\in\ker(\ev)=\Rel\).  The converse follows by
applying \(\ev\) to \(\src z=p+r\).
\end{proof}

This lemma shows that evaluation loses no information in the final
conclusion.  It does not, however, prove that the exceptional pair belongs
to the range.  The invalid argument occurred precisely in the attempted
proof of that membership.

\section{The missing root in the finite Stokes calculation}

Let \(C_0\) be the free abelian group on labelled row occurrences and let
\(C_1\) be the free abelian group on labelled rewrite cells.  If a cell has
one input and a signed list of outputs, define its boundary to be input
minus outputs.  For a finite deterministic rewrite forest with signed root
family \(I\), terminal frontier \(F\), and signed cell sum \(S\), induction
on the processing order gives
\begin{equation}
                         I-F=\partial_{\mathrm{led}}S.       \label{eq:stokes}
\end{equation}
Every internal occurrence cancels, but every root occurrence occurs once
with positive sign.  Cutting away certified subforests changes the frontier
and introduces cut edges; it does not remove an unclassified root.

Let
\[
 \operatorname{Read}(v)=
 \left(
  \Sym^2(\pi)P_2(\operatorname{val}v),
  \ev P_1(\operatorname{val}v)
 \right).
\]
Literal rewrite identities imply
\(\operatorname{Read}(\partial_{\mathrm{led}}S)=0\).  Applying the read to
\eqref{eq:stokes} therefore gives
\begin{equation}
                         \operatorname{Read}(I)
                         =\operatorname{Read}(F),            \label{eq:root-frontier}
\end{equation}
not \(\operatorname{Read}(F)=0\).

In the governing trace, the initial word is the relation-ideal witness
\(\Theta\).  Its terminal factor-two read is zero by the governing weight
condition, but its factor-one read is not zero:
\begin{equation}
 P_1(\Theta)=\widetilde a,
 \qquad
 \ev(\widetilde a)=a.                                      \label{eq:root-read}
\end{equation}
Thus its simultaneous read is ((0,a)).  It cannot be placed among the
zero-read leaves, whole relations, or independently certified correction
families.

The external-boundary formula in the earlier document had the form
\[
 \partial_{\mathrm{led}}S^\circ
    =E-G+Q+Z+\sum_jN_j.
\]
The correct formula must also contain the retained root family
\(I^\circ\), with the sign inherited from \eqref{eq:stokes}.  Consequently
the read only implies
\[
 \operatorname{Read}(E)
   =-\operatorname{Read}(I^\circ)
\]
after the other displayed families have been discharged.  Proving the
right-hand side zero is Step~7 itself.

\section{Exact audit of the canonical exceptional chain}

The preceding diagnosis can be made completely concrete with the objects
already present in the Lean development.  Let
\[
 z_{\mathrm{can}}=
 \texttt{rawCutoffExceptionalRelationLeftChain}
\]
and define its primitive discrepancy
\[
 \Delta_{\mathrm{can}}
   =\src z_{\mathrm{can}}-P_{\mathcal E}.
\]
The boundary theorem already proved in Lean is
\[
                         \partial z_{\mathrm{can}}=B_{\mathcal E}.
\]

For the primitive calculation, abbreviate the five complete primitive
aggregates by
\begin{align*}
 C&=\text{complete factor-two full-label primitive},\\
 T&=\text{raw contextual terminal-two primitive},\\
 N&=\text{non-hole mark-one primitive},\\
 H&=\text{derived-tail hole primitive},\\
 E&=\text{exceptional component primitive}.
\end{align*}
All source-placement errors occurring below are genuine elements of
\(\Rel\).  The already formalized source formulas and the definition of
\(z_{\mathrm{can}}\) give
\begin{equation}
                         \Delta_{\mathrm{can}}
            =-(C+T+N+H+E)+r_0                         \label{eq:canonical-one}
\end{equation}
for some \(r_0\in\Rel\).  No component is asserted to be a relation in
this calculation: all relation terms are the explicitly constructed
aggregate placement errors.

The verified mark-one splits give
\begin{equation}
                         N+H+E=P_{\mathrm{trace}},           \label{eq:mark-split}
\end{equation}
and the verified raw full-label Stokes equation gives
\begin{equation}
 C+P_{\mathrm{trace}}+T
       =P_1(\Theta)-r_1                                    \label{eq:raw-stokes}
\end{equation}
for the grouped terminal one-factor relation \(r_1\in\Rel\).  Substitution
in \eqref{eq:canonical-one} yields the exact identity
\begin{equation}
                         \Delta_{\mathrm{can}}
             =-P_1(\Theta)+r_{\mathrm{can}}                 \label{eq:canonical-root}
\end{equation}
with \(r_{\mathrm{can}}\in\Rel\).  Applying evaluation and
\eqref{eq:root-read} proves the following.

\begin{theorem}[canonical root audit]\label{thm:canonical-audit}
\begin{equation}
                         \ev(\Delta_{\mathrm{can}})=-a.      \label{eq:canonical-eval}
\end{equation}
\end{theorem}

This theorem was checked in Lean, without any new axiom or placeholder, as
the statement
\[
\begin{split}
&\texttt{evaluation n L data}\;\bigl(
  \texttt{rawCutoffExceptionalRelationLeftDiscrepancy}\;\cdots\\
&\hspace{35mm}
  (\texttt{rawCutoffExceptionalRelationLeftChain}\;\cdots)\bigr)
  =-a.
\end{split}
\]
Its proof uses only the existing primitive formulas and the existing three
mark-one splits, together with the two existing theorems
\begin{gather*}
 \texttt{rawPBWPrimitiveTheta\_eq\_complete\_add\_cutoff},\\
 \texttt{pbwPrimitive\_rawCompleteCutoff\_external},
\end{gather*}
and the governing identities \texttt{pbwPrimitive\_theta} and
\texttt{evaluates}.

\begin{corollary}\label{cor:circular}
The assertion \(\ev(\Delta_{\mathrm{can}})=0\) is equivalent to \(a=0\).
In particular it is not an independent lemma from which Step~7 may be
deduced.
\end{corollary}

\begin{proof}
By Theorem~\ref{thm:canonical-audit}, the asserted equality is
\(-a=0\), which is equivalent to \(a=0\).
\end{proof}

This is the exact algebraic form of the circular cancellation observed in
the earlier formalization attempts.

\section{Why the terminal-package proposition does not repair the gap}

The earlier document introduced an aggregate terminal chain
\(z_{\mathrm{ter}}\) and asserted, after a descending continuation, that
its boundary and evaluated source defect were exactly the retained terminal
edges and the intermediate diagonal differences.  That assertion was not a
consequence of the stated local rules.

There are two separate problems.
\begin{enumerate}
\item The continuation was defined only verbally.  The existing normalized
component frontier records component descendants, whereas a terminal source
tensor requires a genuine full contextual relation.  The theorem
\texttt{componentPBWFrontier\_factorTwo\_context\_ne\_hole} supplies such a
relation for each supported factor-two state, but replacing the component
by the full prefix creates the preceding-prefix terms already isolated in
\texttt{ContextualTerminalDiagonalAudit}.  The available theorem proves
their aggregate boundary, not the matching primitive formula.
\item Even if that new continuation were constructed, its Stokes boundary
would still contain the root family described in Section~4.  Omitting the
root makes the claimed terminal-package identity false; retaining it leaves
the governing term \(-P_1(\Theta)\), exactly as in
\eqref{eq:canonical-root}.
\end{enumerate}

Thus Proposition~5.4 of the earlier version was not a smaller lemma.  It
contained both the old missing primitive assertion and an omitted root
term.  Attempting to formalize it would very likely be more complicated
than the previous attempt, because it would additionally require a new
state continuation and its preservation theorems.

The evaluated cokernel does not alter this conclusion.  A quotient can
package simultaneous boundary and primitive equations, but it cannot turn
the nonzero root class into zero.  By Lemma~\ref{lem:range}, proving that
the exceptional class vanishes is exactly proving the desired lift.

\section{The correct root-inclusive target}

A noncircular Step~7 proof must keep the root contribution and identify it
with the terminal coordinate of a genuine terminal source cycle.  The
minimal coordinate-level form is the following.

\begin{proposition}[root-inclusive terminal certificate]\label{prop:correct-target}
For the governing witness there are a genuine cycle
\[
 c\in Z_1K_2(D_n\hookrightarrow A_n)
\]
and a terminal preimage \(s\) whose underlying free-model element is
\(\src(c)\), such that
\begin{equation}
 \operatorname{terminalEval}(s)
       =\operatorname{coord}_C(a).                          \label{eq:root-certificate}
\end{equation}
\end{proposition}

Point~6, already formalized, annihilates the terminal evaluation of the
source primitive of every genuine cycle.  Hence
\eqref{eq:root-certificate} implies
\(\operatorname{coord}_C(a)=0\), and injectivity of the cyclic coordinate
gives \(a=0\).

This target matches the existing consumer
\[
             \texttt{eq\_zero\_of\_terminalSourceCoordinate}.
\]
Equivalently, one may construct a
\texttt{TerminalQuadraticPacket} and apply the existing theorem
\[
                    \texttt{eq\_zero\_of\_terminalPacket}.
\]
Either route must establish a root-inclusive trace certificate.  It must
not first discard the root and try to prove the exceptional edge zero in
isolation.

The original manuscript's closed-square equality is one mathematical proof
of Proposition~\ref{prop:correct-target}: the terminal root, every
horizontal and vertical intermediate diagonal, the oriented factor-two
edge, and the terminal one-factor relations are read in one signed ledger.
For Lean, however, the missing work remains the construction of that one
cycle/certificate from the present collectors.  The current library has
already formalized the local PBW identities, the diagonal transgression,
the terminal realization, and the final consumer; it has not yet formalized
the required root-inclusive certificate.

\section{Implementation consequence}

There is presently no basis for promising that Step~7 will be easier than
the previous attempt.  The safe conclusion of this audit is narrower:
\begin{itemize}
\item do not implement the evaluated cokernel or the proposed terminal
package;
\item retain the current \texttt{Finsupp} collectors and aggregate source
formulas;
\item make the root term explicit in every invariant; and
\item target the existing coordinate-level terminal consumer rather than a
stronger full-relation lift of an isolated edge.
\end{itemize}

This avoids known false goals and unnecessary infrastructure.  It does not
remove the one genuinely hard theorem: the root-inclusive terminal
certificate.

\section{Four-point implementation plan for all of Step~7}

The first two points contain all mathematical and bookkeeping complexity.
The last two are deliberately short applications of existing code.

\begin{enumerate}
\item \textbf{Freeze and verify the root-inclusive invariant.}
  Add Theorem~\ref{thm:canonical-audit} to a small audit file in the Lean
  library and prove the corresponding identity for every proposed partial
  continuation.  State the final collector invariant with an explicit root
  summand; never use an external-boundary partition lacking that summand.
  Fix the final interface to be Proposition~\ref{prop:correct-target}, or
  the equivalent \texttt{TerminalQuadraticPacket} certificate.  This point
  prevents another long attempt at a goal which is secretly equivalent to
  \(a=0\).

\item \textbf{Construct the single root-inclusive terminal certificate.}
  Work on the existing complete provenance and normalized-component
  \texttt{Finsupp} ledgers.  Retain the initial governing rows, the paired
  intermediate diagonals, the preceding-prefix corrections, and the
  factor-two terminal rows in the same aggregate.  Prove simultaneously
  that the resulting terminal chain is closed and that its complete source
  primitive has the governing terminal coordinate.  Use the existing
  \texttt{Assembly}/\texttt{HigherCharacters} transgression for paired
  diagonals, the existing full-relation source-placement formulas for
  factor-two rows, and the existing comb support results for the one-factor
  frontier.  This is the only new substantive theorem.  If its proof begins
  to require a second independent primitive collector, stop and strengthen
  the common ledger invariant instead of splitting the two reads.

\item \textbf{Invoke the existing terminal consumer.}
  Package the cycle and coordinate equation from Point~2 and apply
  \texttt{eq\_zero\_of\_terminalSourceCoordinate}; alternatively construct
  the already defined \texttt{TerminalQuadraticPacket} and apply
  \texttt{eq\_zero\_of\_terminalPacket}.  Point~6 and injectivity of the
  cyclic coordinate then give \(a=0\).  No new PBW calculation belongs in
  this point.

\item \textbf{Export and test Step~7.}
  Export the unconditional capstone expected by \texttt{StepSeven} and
  \texttt{Main}, specialize the existing reduction theorem, and run the
  complete project build and low-class regression tests, including the
  empty intermediate-diagonal case \(n=2\).  This point is module wiring and
  testing only.
\end{enumerate}

\section{Conclusion}

The double-check changes the conclusion decisively.  The evaluated range
criterion is correct, but the proposed proof of exceptional range membership
omitted the initial governing root.  The exact Lean-checked identity
\(\ev(\Delta_{\mathrm{can}})=-a\) shows that the omission concealed the
whole Step~7 conclusion.  The previous proof is therefore withdrawn, and
the PDF should be used as an implementation warning and corrected target,
not as a proof of the exceptional lifting lemma.

\end{document}
